

\frame{
\ftitle{Vorwort}{\black}
	
	Diese Slides sollen gute Übungsaufgaben für das Fach ''Numerisches Programmieren'' bieten.
	Ich versuche hier, zu den wichtigsten Themen Beispielaufgaben und weitere Aufgaben selber zu erstellen,
	sodass ihr mit den Slides direkt lernen könnt.
	
	Ich wünsche viel Erfolg beim lernen,

	Hendrik
}

\frame{
	\ftitle{Bitte}{\black}
	
	Wenn ihr Fehler oder Ungereimheiten findet, schreibt mir bitte direkt eine E-mail.\\
	Adresse: \fat{\npauthormail{}}
	
	Ich würde darum bitten, dass im Betreff ''Tutorium'' vorkommt, \\
	dadurch wird der Mail eine höhere Priorität zugeordnet.
	
}

\frame{
	\ftitle{Disclaimer}{\red}
	
	Diese Zusammenfassung wurde von einem Tutor erstellt und erhebt keinen Anspruch auf	\fat{Lesbarkeit, Verständlichkeit, Richtigkeit oder Vollständigkeit}.

	Sie soll lediglich als \fat{Lernunterstützung und Vorbereitungsmöglichkeit} dienen.
}

\frame{
	\ftitle{Nutzung}{\blue}
	
	\begin{itemize}
		\item Grüne Slides: Kapitelfolien
		\item Orangene Slides: Kündigen ein Thema des Kapitels an
		\item Graue Slides: Beschreibt Aufgaben
		\item Gelbe Slides: Lösungen zu vorheriger Aufgaben Slide
	\end{itemize}
	Ich würde empfehlen, mir das Thema bei Orangenen Slides anzuschauen. Dann zur nächsten, grauen Folie zu wechseln und direkt alle Aufgaben dieser Folie zu bearbeiten. \\
	Erst wenn man mit allen Aufgaben einer Folie fertig ist, sollte zu der Lösungsfolie gewechselt werden.
	\n 
	Wichtig: Es kommen nie zwei Aufgaben Slides hintereinander, sondern nach einer kommt immer die dazugehörige Slide mit den Lösungen.
}